% ============================================================================
% Section 5.6 — Graphing Solutions to Inequalities on a Number Line
% CCSS 7.EE.B.4.B
% ============================================================================
\section{Graphing Inequalities on a Number Line}

\begin{stepGoal}
\begin{itemize}
    \item Graph inequality solutions using open and closed circles
    \item Interpret graphs in the context of real-world problems
\end{itemize}
\end{stepGoal}

\begin{stepVocabBox}
    \vocabItem{Open Circle}{A hollow dot ($\circ$) showing the endpoint is \textbf{not} included ($<$ or $>$).}
    \vocabItem{Closed Circle}{A filled dot ($\bullet$) showing the endpoint \textbf{is} included ($\leq$ or $\geq$).}
\end{stepVocabBox}

\begin{stepsCard}[How to Graph an Inequality on a Number Line]
    \stepItem{\textbf{Solve} the inequality to get the variable alone (e.g., $x < 3$).}
    \stepItem{Draw a \textbf{circle} at the boundary number: \textbf{open} for $<$ or $>$, \textbf{closed} for $\leq$ or $\geq$.}
    \stepItem{\textbf{Shade} the direction of all solutions: left for ``less than,'' right for ``greater than.''}
    \stepItem{\textbf{Interpret}: explain what the graph means in the problem's context.}
\end{stepsCard}

% --- Worked Example 1 ---
\begin{stepExample}{Solve $2x - 1 < 5$ and describe the graph.}
    \stepShow{1} Add $1$: $2x < 6$. Divide by $2$: $x < 3$.

    \stepShow{2} Open circle at $3$ (not included).

    \stepShow{3} Shade to the left of $3$.

    \stepShow{4} All numbers less than $3$ are solutions.
    \stepResult{$x < 3$; open circle at $3$, shade left.}
\end{stepExample}

% --- Worked Example 2 ---
\begin{stepExample}{Solve $-2x \leq 8$ and describe the graph.}
    \stepShow{1} Divide by $-2$ and flip the sign: $x \geq -4$.

    \stepShow{2} Closed circle at $-4$ (included).

    \stepShow{3} Shade to the right of $-4$.

    \stepShow{4} All numbers greater than or equal to $-4$ are solutions.
    \stepResult{$x \geq -4$; closed circle at $-4$, shade right.}
\end{stepExample}

\mascotStepTip{Think of the circle as a gate. Open means ``you can't come in.'' Closed means ``you're included!''}

% ============================================================================
% PRACTICE
% ============================================================================
\newpage
\begin{stepPractice}{Graphing Inequalities Practice}
    \resetProblems

    \stepPracticeHeader[step1Color]{Open or Closed?}
    \begin{multicols}{2}
    \prob $x > 5$: open or closed circle? \answerBlank[2cm]
    \answer{open}
    \prob $x \leq -2$: open or closed circle? \answerBlank[2cm]
    \answer{closed}
    \end{multicols}

    \stepPracticeHeader[step2Color]{Solve and Describe the Graph}
    \prob Solve $x + 4 \geq 7$. Describe the graph. \answerBlank[4cm]
    \answer{$x \geq 3$; closed circle at $3$, shade right}

    \prob Solve $3n - 6 > 0$. Describe the graph. \answerBlank[4cm]
    \answer{$n > 2$; open circle at $2$, shade right}

    \stepPracticeHeader[step3Color]{Real-World Graphing}
    \trueOrFalse{The graph of $x < 5$ uses a closed circle at $5$.}
    \answerTF{False}

    \wordProblem{A parking lot charges $\$3$ per hour. You have $\$20$. Write and solve an inequality for the hours $h$ you can park, and describe the graph.}{hours}
    \answer{$3h \leq 20$; $h \leq 6\frac{2}{3}$; closed circle at $6\frac{2}{3}$, shade left}
\end{stepPractice}
